\section{DB間の移行等}

\subsection{規約型DBの統合}

出題例なし

\newpage

\subsection{規約型DBの分割}

\begin{InyouJobunbox}
  \begin{itemize}
    \item 法第75条
    \item 令第51条
    \item 則第87条の2、第91条
    \item 特例的扱い第7、第10
  \end{itemize}
\end{InyouJobunbox}


\begin{fillblank}[年金2 2023 問題1(2)　\&　数理人試験 H29 問題1 設問1]
  \noindent \textbf{◯DB則}
  \begin{lawstructure}
    \begin{lawarticle}{第八十七条の二}{確定給付企業年金の分割時に移換する積立金の額の算定方法}
      法第七十五条第一項の規定により規約型企業年金を分割する場合又は法第七十七条第一項の規定により基金を分割する場合における分割された規約型企業年金の資産管理運用機関又は分割により設立された基金（以下この項において「移換先確定給付企業年金」という。）に移換する積立金の額の算定方法は、次の各号のいずれかの方法とする。

      \lawitem{一}
      当該分割を行う日（以下この号において「分割日」という。）の前日における
      当該分割を行う規約型企業年金又は基金の積立金（以下この項において「分割
      時積立金」という。）の額を（　　A　　）、直近の財政計算の計算基準日、そ
      の前の財政計算の計算基準日又は（　　B　　）
      における次に掲げる額のいずれかに応じて按分する方法

      \hangedlawitem{イ}{3}
      （　　C　　）

      \hangedlawitem{ロ}{3}
      数理債務の額

      \hangedlawitem{ハ}{3}
      数理債務の額から特別掛金額の予想額の現価と第四十七条に定める掛金の額の予想額の現価を合算した額を控除した額

      \hangedlawitem{ニ}{3}
      （　　A　　）、直近の財政計算の計算基準日若しくはその前の財政計算の計算
      基準日を法第六十条第三項に規定する事業年度の末日とみなして同項の規定に基
      づき算定した最低積立基準額又は（　　B　　）に
      おける最低積立基準額

      \lawitem{二}
      次に定める額のうち、移換先確定給付企業年金に係る額の合計額とする方法
      （分割時積立金の額が本号イの算定に用いる前号に掲げる額を（　　D　　）場合に限る。）
      
      \hangedlawitem{イ}{3}
      前号に掲げるいずれかの額のうち（　　E　　）に係る部分の額（分割時積立
      金の額が前号に掲げるいずれかの額のうち（　　E　　）に係る部分の額の合計額
      を下回る場合にあっては、当該分割時積立金の額を当該前号に掲げるいずれか
      の額のうち（　　E　　）に係る部分の額に応じて按分して得た額）

      \hangedlawitem{ロ}{3}
      分割時積立金の額からイに掲げる額の合計額を控除した額につき、
      本号イの算定に用いる前号に掲げる額のうち（　　F　　）（（　　E　　）を除く。）に係る部分の額に応じて按分して得た額

      \lawitem{三}
      （　　G　　）しないよう移換先確定給付企業年金に移換する積立金の額を定める方法（（　　H　　）の場合において、分割により（　　G　　）することが見込まれる場合に限る。）

      \lawitem{四}
      その他厚生労働大臣が定める方法（厚生労働大臣が定める場合に限る。）

      \lawparagraph{2}
      （略）

    \end{lawarticle}
  \end{lawstructure}
\end{fillblank}

\begin{answerbox}
  \begin{enumerate}[A: ]
    \item 分割日の前日
    \item 分割日が属する事業年度の前事業年度の末日
    \item 通常予測給付額の現価
    \item 下回る
    \item 受給権者等
    \item 加入者
    \item 積立割合、調整率又は超過比率が減少
    \item リスク分担型企業年金
  \end{enumerate}
\end{answerbox}

\begin{shortanswer}[年金1 2021 問題2(1)\textcircled{4}]
  確定給付企業年金法施行規則第87 条の2 第1 項第3 号に定める、リスク分担型企業年金が定
めることができる規約型企業年金または基金の分割時に移換する積立金の額の算定方法につい
て簡記しなさい。なお、当該算定方法を選択する場合の条件（制約）についても明記するこ
と。
\end{shortanswer}

\begin{answerbox}
  ・積立割合、調整率又は超過比率が減少しないよう移換先確定給付企業年金に移換する積立金の額
を定める方法。なお、分割により積立割合、調整率又は超過比率が減少すると見込まれる場合に限
る。

・なお、積立割合は調整前給付現価相当額と給付財源の比率を、超過比率は調整前給付現価相当額に
対する給付財源から調整前給付現価相当額と財政悪化リスク相当額の２分の１の額を合算した額
を控除した額の比率をそれぞれ指す。
\end{answerbox}

\newpage

\subsection{基金の合併}

\begin{InyouJobunbox}
  \begin{itemize}
    \item 法第76条
    \item 令第52条、第53条の2
    \item 則第92条
  \end{itemize}
\end{InyouJobunbox}


\begin{fillblank}[年金1 H29 問題1(1)B]
  \noindent \textbf{◯DB法}
  \begin{lawstructure}
    \begin{lawarticle}{第七十六条}{基金の合併}
      基金は、合併しようとするときは、厚生労働大臣の認可を受けなければならない。

      \lawparagraph{2}
      前項の認可の申請は、代議員会における（　　A　　）の多数による議決を経て行わなければならない。
    
      （以下略）
    \end{lawarticle}
  \end{lawstructure}
\end{fillblank}

\begin{answerbox}
  \begin{enumerate}[A: ]
    \item 代議員の定数の四分の三以上
  \end{enumerate}
\end{answerbox}

\begin{explanationbox}
  類題：年金１ 2023 問題1(3)ア
\end{explanationbox}

\begin{truefalse}[年金2 2022 問題1(1)(エ)\ajRoman{1}]
  以下の文章の正誤を判定し、誤っている場合はその内容を改めなさい。
\\

  基金の合併にかかる認可の申請は、代議員会における代議員の定数の四分の三以上の多
数による議決を経て行わなければならない。
\end{truefalse}

\begin{answerbox}
  \textbf{正}
\end{answerbox}

\begin{truefalse}[年金2 2022 問題1(1)(エ)\ajRoman{4}]
  以下の文章の正誤を判定し、誤っている場合はその内容を改めなさい。
\\

  合併又は分割により設立された基金は、基金の設立による公告に併せて、合併により消
滅した基金又は分割前の基金の名称及び所在地を公告しなければならない。
\end{truefalse}

\begin{answerbox}
  \textbf{正}
\end{answerbox}

\newpage

\subsection{基金の分割}

\begin{InyouJobunbox}
  \begin{itemize}
    \item 法第77条
    \item 令第52条、第53条の2
    \item 則第87条の2、第93条
    \item 承認認可基準別添
  \end{itemize}
\end{InyouJobunbox}

\begin{fillblank}[年金1 H29 問題1(1)C、D、E]
  \noindent \textbf{◯DB法}
  \begin{lawstructure}
    \begin{lawarticle}{第七十七条}{基金の分割}
      基金は、分割しようとするときは、厚生労働大臣の認可を受けなければならない。

      \lawparagraph{2}
      基金の分割は、（　　A　　）について行うことはできない。

      \lawparagraph{3}
      分割を行う場合においては、分割により設立される基金の加入者となるべき厚生年金保険の被保険者又は分割後存続する基金の加入者である厚生年金保険の被保険者の数が、第十二条第一項第四号（基金を共同して設立している場合にあっては、同項第五号）の政令で定める数以上であるか、又は当該数以上となることが見込まれなければならない。

      \lawparagraph{4}
      分割によって基金を設立するには、分割により設立される基金の実施事業所となるべき厚生年金適用事業所の事業主が規約を作り、その他設立に必要な行為をしなければならない。

      \lawparagraph{5}
      分割により設立された基金は、分割により消滅した基金又は分割後存続する基金の（　　B　　）を承継する。

      \lawparagraph{6}
      前項の規定により承継する（　　C　　）は、厚生労働大臣の認可を受けなければならない。

      \lawparagraph{7}
      前条第二項の規定は、第一項及び前項の認可の申請を行う場合について準用する。
    \end{lawarticle}
  \end{lawstructure}
\end{fillblank}

\begin{answerbox}
  \begin{enumerate}[A: ]
    \item 実施事業所の一部
    \item 権利義務の一部
    \item 権利義務の限度
  \end{enumerate}
\end{answerbox}

\begin{truefalse}[年金2 2022 問題1(1)(エ)\ajRoman{2}]
    以下の文章の正誤を判定し、誤っている場合はその内容を改めなさい。
\\

  基金の分割は、実施事業所の一部について行うことはできない。
\end{truefalse}

\begin{answerbox}
  \textbf{正}
\end{answerbox}


\begin{truefalse}[年金2 2022 問題1(1)(エ)\ajRoman{3}]
    以下の文章の正誤を判定し、誤っている場合はその内容を改めなさい。
\\

  基金の分割を行う場合においては、分割により設立される基金の加入者となるべき厚生
年金保険の被保険者又は分割後存続する基金の加入者である厚生年金保険の被保険者の
数が三百人以上でなければならない。
\end{truefalse}


\begin{answerbox}
  \textbf{誤}

正しくは「300人以上であるか、300人以上となることが見込まれていなければならない」  
\end{answerbox}

\newpage

\subsection{実施事業所の増加または減少}


\begin{InyouJobunbox}
  \begin{itemize}
    \item 法第78条
    \item 則第88条〜第88条の2
    \item 解釈5-3
    \item 特例的扱い第7、第8
  \end{itemize}
\end{InyouJobunbox}

\begin{fillblank}[数理人試験 2020 問題1 設問2　\&　年金2 H28 問題1(3)C〜I]
  \noindent \textbf{◯DB則}
  \begin{lawstructure}
    \begin{lawarticle}{第八十八条}{実施事業所の減少に係る掛金の一括徴収}
      法第七十八条第三項の厚生労働省令で定める事由は、次のとおりとする。

      \lawitem{一}
      実施事業所の事業主が、分割又は事業の譲渡により他の実施事業所の事業主以外の事業主にその事業の全部又は一部を承継させる場合

      \lawitem{二}
      前号に規定する場合のほか、規約で定めるところにより、実施事業所に使用される当該確定給付企業年金の（　　A　　）場合
    \end{lawarticle}


    \begin{lawarticle}{第八十八条の二}{}
      法第七十八条第三項の厚生労働省令で定める計算方法は、次のいずれかの方法とする。

      \lawitem{一}
      当該減少に係る実施事業所（以下この条において「減少実施事業所」という。）が（　　B　　）としたならば減少実施事業所の事業主が拠出することとなる（　　C　　）とする方法

      \lawitem{二}
      前号の方法により計算した額に（　　D　　）で定めるところにより次のイからハまでに掲げる場合の区分に応じ、当該イからハまでに定める額を加算した額とする方法

      \hangedlawitem{イ}{3}
      減少実施事業所が減少する日（以下この条において「減少日」という。）において、積立金の額が当該減少日を法第六十条第二項に規定する事業年度の末日とみなして同項の規定に基づき算定した（　　E　　）を下回ることが見込まれる場合　当該下回る額の見込額を償却するために必要となる掛金の額のうち減少実施事業所が（　　B　　）としたならば減少実施事業所の事業主が拠出することとなることが見込まれる掛金の額として合理的に計算した額

      \hangedlawitem{ロ}{3}
      減少日において、時価により評価した積立金の額が前回の財政計算の計算基準日において用いた第四十八条第一項に規定する方法で評価した積立金の額を下回ることが見込まれる場合　当該下回る額の見込額を償却するために必要な掛金の額のうち減少実施事業所が（　　B　　）としたならば減少実施事業所の事業主が拠出することとなることが見込まれる掛金の額として合理的に計算した額

      \hangedlawitem{ハ}{3}
      減少実施事業所の減少に併せて掛金の額の（　　F　　）をするとした場合において、イ又はロ以外の要因により掛金の額が増加することとなる場合　当該イ又はロ以外の要因により増加することとなる掛金の額のうち減少実施事業所の事業主が拠出すべき額として合理的に計算した額

      \lawitem{三}
      減少日における積立金の額が、当該日を法第六十条第三項に規定する事業年度の末日とみなして同項の規定に基づき算定した（　　G　　）を下回ることが見込まれる場合において、当該下回る額の見込額のうち減少実施事業所に係る分として規約で定めるところにより合理的に計算した額とする方法

      \lawitem{四}
      （　　H　　）とする方法

      \lawitem{五}
      （　　I　　）とする方法

      \lawitem{六}
      その他厚生労働大臣が定めるところにより計算した額とする方法（第八十七条の二第一項第四号の厚生労働大臣が定める場合に限る。）

      \lawparagraph{2}
      前項第一号の（　　C　　）の計算に用いる予定利率は、第四十三条第二項第一号の規定に基づき定めた予定利率とする。

      \lawparagraph{3}
      事業主等は、規約で定めるところにより、第一項に規定する方法で計算した額に、減少実施事業所が減少しないとしたならば減少実施事業所の事業主が負担することとなる第四十五条第一項に規定する（　　J　　）の額を加算することができる。
    \end{lawarticle}
  \end{lawstructure}
\end{fillblank}


\begin{answerbox}
  \begin{enumerate}[A: ]
    \item 加入者の数が減少する
    \item 減少しない
    \item 特別掛金額の予想額の現価
    \item 規約
    \item 責任準備金の額
    \item 再計算
    \item 最低積立基準額
    \item 第一号または第三号のうちいずれか大きい額
    \item 第二号または第三号のうちいずれか大きい額
    \item その他の掛金の額
  \end{enumerate}
\end{answerbox}

\begin{truefalse}[年金1 2025 問題1(3)ア]
  下線部の正誤を判断し、誤っている場合はその内容を改めなさい。
  \\

  事業主等がその実施事業所を増加させ、又は減少させようとするときは、その増加又は減少に
係る厚生年金適用事業所の\underline{労働組合等（当該厚生年金適用事業所に使用される厚生年金保険の
被保険者の過半数で組織する労働組合があ}\\
\underline{るときは当該労働組合、当該厚生年金保険の被保険
者の過半数で組織する労働組合がないときは当該厚生年金保険の}\\
\underline{被保険者の過半数を代表す
る者）の４分の３以上の同意}を得なければならない。
\end{truefalse}

\begin{answerbox}
  \textbf{誤}

  事業主の全部の同意及び労働組合等（当該厚生年金適用事業所に使用される厚生年金
保険の被保険者の過半数で組織する労働組合があるときは当該労働組合、当該厚生年
金保険の被保険者の過半数で組織する労働組合がないときは当該厚生年金保険の被保
険者の過半数を代表する者）の全部の同意
\end{answerbox}

\newpage

\subsection{実施事業所の減少の特例}


\begin{InyouJobunbox}
  \begin{itemize}
    \item 法第78条の2
    \item 令第48条の2
    \item 則第88条の3
  \end{itemize}
\end{InyouJobunbox}

\begin{fillblank}[年金1 2021 問題1(1)C　\&　年金2 H28 問題1(3)A〜C]
  \noindent \textbf{◯DB法}
  \begin{lawstructure}
    \begin{lawarticle}{第七十八条の二}{確定給付企業年金を実施している事業主が二以上である場合等の実施事業所の減少の特例}
      確定給付企業年金を実施している事業主が二以上である場合又は基金が二以上の事業主により設立された場合において、事業主等が一の事業主の実施事業所の全てを減少させようとする場合であって次に掲げる要件を満たすときは、前条第一項の規定にかかわらず、厚生労働大臣の承認（確定給付企業年金が基金型企業年金である場合にあっては、認可）を受けて、当該実施事業所を減少させることができる。

      \lawitem{一}
      減少させようとする実施事業所の事業主が確定給付企業年金を（　　A　　）であると認められること。

      \lawitem{二}
      基金の場合にあっては、基金の（　　B　　）の数が、当該実施事業所を減少させた後においても、第十二条第一項第四号（基金を共同して設立している場合にあっては、同項第五号）の政令で定める数以上であるか、又は当該数以上となることが見込まれること。

      \lawitem{三}
      当該実施事業所の減少に伴い他の実施事業所の事業主の掛金が増加することとなる場合にあっては、（　　C　　）において、当該減少に係る実施事業所の事業主が、当該増加する額に相当する額として厚生労働省令で定める計算方法のうち当該（　　C　　）で定めるものにより算定した額を、掛金として一括して拠出する旨を定めていること。
    \end{lawarticle}
  \end{lawstructure}
\end{fillblank}


\begin{answerbox}
  \begin{enumerate}[A: ]
    \item 継続することが困難
    \item 加入者
    \item 規約
  \end{enumerate}
\end{answerbox}


\begin{fillblank}[年金1 2021 問題1(1)C、D　\&　年金2 H28 問題1(3)A,C,J.K]
  \noindent \textbf{◯DB則}
  \begin{lawstructure}
    \begin{lawarticle}{第八十八条の三}{実施事業所の減少の特例を適用する場合の手続等}
      法第七十八条の二第一号の確定給付企業年金を（　　A　　）であると認められることは、同条の規定による実施事業所の減少に関する事項を（　　B　　）に定めた場合であって、当該事項を（　　B　　）に定めた日以後に減少させようとする実施事業所の事業主が（　　C　　）（当該事業主がその責に帰することができない事由により掛金を納付することができない期間がある場合にあっては、当該期間に係る掛金額に相当する額を除く。）を超えて掛金の納付を怠ったこととする。

      \lawparagraph{2}
      事業主等は、法第七十八条の二の規定により実施事業所を減少させようとする場合には、当該実施事業所の事業主に対し、掛金の納付を怠った理由について（　　D　　）を与えなければならない。

    （以下略）
    \end{lawarticle}
  \end{lawstructure}
\end{fillblank}

\begin{answerbox}
  \begin{enumerate}[A: ]
    \item 継続することが困難
    \item 規約
    \item 一年分に相当する額
    \item 弁明の機会
  \end{enumerate}
\end{answerbox}

\newpage


\subsection{給付の支給に関する権利義務の他のDBへの移転}


\begin{InyouJobunbox}
  \begin{itemize}
    \item 法第79条
    \item 令第49条〜第50条、第52条〜第53条
    \item 則第87条の2、第89条〜第89条の2、第94条
    \item 解釈第5-4
    \item 特例的扱い第10
  \end{itemize}
\end{InyouJobunbox}

\begin{shortanswer}[数理人試験 H23 問題2(1)]
  規約型企業年金の実施事業所から他の規約型企業年金の実施事業所へ権利義務の移転を行
  うとき、事業主が移転の承認の申請をするために、移転される確定給付企業年金の加入者
  に関して得なければならない同意手続きの内容について簡記せよ。
\end{shortanswer}

\begin{answerbox}
  \begin{itemize}
    \item 実施事業所に使用される移転加入者の過半数で組織する労働組合があるときは当該労働
 組合、当該移転加入者の過半数で組織する労働組合がないときは当該移転加入者の過半
 数を代表する者の同意を得る必要がある。
    \item 同意は、移転加入者が使用される実施事業所が二以上であるときは、各実施事業所につ
 いて得なければならない。
    \item 尚、移転加入者とは権利義務が移転される移転確定給付企業年金の加入者のことを指す。
  \end{itemize}
\end{answerbox}

\begin{shortanswer}[数理人試験 H23 問題2(2)]
  基金型企業年金の実施事業所から他の基金型企業年金の実施事業所へ権利義務の移転を行
  うとき、基金が移転の認可の申請をするために、上記(1)と同様の同意手続きに加えて、
  代議員に関して得なければならない同意手続きの内容について簡記せよ。
\end{shortanswer}

\begin{answerbox}
  移転確定給付企業年金基金の代議員のうち、移転加入者以外の加入者が使用される実施事
業所に係る代議員の定数の四分の三以上の同意を得る必要がある。
\end{answerbox}

\newpage

\subsection{規約型DB $\leftrightarrow$ 基金型DB}


\begin{InyouJobunbox}
  \begin{itemize}
    \item 法第80条〜第81条
    \item 令第52条〜第53条、第58条
    \item 則第95条〜第96条
  \end{itemize}
\end{InyouJobunbox}

\begin{shortanswer}[年金1 2019 問題2(3)\textcircled{2}]
  次の A・B の 2 つの確定給付企業年金間の移行を行う場合において、承認（認可）申請を行
うにあたって移転側・承継側で必要な手続きについてそれぞれ簡記しなさい。
\begin{enumerate}[A]
  \item 基金型企業年金から規約型企業年金へ移行するとき
  \item 規約型企業年金から基金型企業年金へ移行するとき
\end{enumerate}
ただし、いずれにおいても実施事業所は 1 つであり、承継側制度を新たに設立するものとし、
移行に際して給付設計の変更などは行わないものとする。
\end{shortanswer}

\begin{answerbox}
  \begin{enumerate}[A]
    \item ＜移転側（基金型）＞
    
代議員会における代議員の定数の四分の三以上の多数により議決すること

＜承継側（規約型）＞

厚生年金保険の被保険者の過半数で組織する労働組合があるときは当該労働組合、当該労働組合がないときは被保険者の過半数を代表する者（以下、労働組合等）から、規約型企業年金を実施することに係る同意および給付に関する権利義務を承継することに係る同意を得ること
\\

    \item ＜移転側（規約型）＞
    
労働組合等から、基金型企業年金へ給付に関する権利義務を移転することに係る同意を得ること、企業年金基金を設立しようとする厚生年金適用事業所の事業主に対して、給付に関する権利義務の移転を申し出ること

＜承継側（基金型）＞

労働組合等から、基金型企業年金を実施することに係る同意及び当該基金が給付に関する権利義務を承継することに係る同意を得ること
  \end{enumerate}
\end{answerbox}

\newpage

\subsection{他のDBへの脱退一時金相当額の移換}


\begin{InyouJobunbox}
  \begin{itemize}
    \item 法第81条の2
    \item 令第50条の2〜第50条の4
    \item 則第32条の2〜第32条の3、第89条の3〜第89条の6
  \end{itemize}
\end{InyouJobunbox}

\begin{fillblank}[年金1 2020 問題1(1)E]
  \noindent \textbf{◯DB則}
  \begin{lawstructure}
    \begin{lawarticle}{第八十九条の四}{脱退一時金相当額の算定の基礎となった期間の一部を合算する場合における算定方法}
      令第五十条の三の規定により脱退一時金相当額の算定の基礎となった期間の一部を、当該中途脱退者に係る移換先確定給付企業年金の加入者期間に算入するときは、次の各号に掲げる要件を満たす算定方法によらなければならない。

      \lawitem{一}
      移換先確定給付企業年金の規約に照らして当該移換された脱退一時金相当額の算定の基礎となる期間を算定すること。ただし、算定された期間が脱退一時金相当額の算定の基礎となった期間を超える場合にあっては、当該算定の基礎となった期間とすること。

      \lawitem{二}
      脱退一時金相当額の算定の基礎となった期間を算入しないこととする場合にあっては、移換先確定給付企業年金の加入者であった期間が（　　A　　）である者に限り、その旨を規約で定めること。

      \lawitem{三}
      その他当該中途脱退者について不当に差別的なものでなく合理的な計算方法であると認められること。
    \end{lawarticle}
  \end{lawstructure}
\end{fillblank}

\begin{answerbox}
  \begin{enumerate}[A: ]
    \item 一年未満
  \end{enumerate}
\end{answerbox}

\begin{truefalse}[年金1 2018 問題1(1)\textcircled{4}]
  下線部の正誤を判断し、誤っている場合はその内容を改めなさい。
  \\

  リスク分担型企業年金でない確定給付企業年金の場合、資産管理運用機関等が脱退一時金
相当額等の移換を受けた者に事業主等が支給する一時金（年金として支給する老齢給付金の
支給を開始した後に支給する一時金を除く）の額は、規約で定める方法により計算した額と
移換を受けた脱退一時金相当額等の額の\underline{合計額}とすることとされている。
\end{truefalse}

\begin{answerbox}
  \textbf{誤}

  いずれか高い額
\end{answerbox}

%この問題、本当にこの位置で合ってる？
\begin{truefalse}[年金1 H29 問題1(5)\textcircled{5}]
    下線部の正誤を判断し、誤っている場合はその内容を改めなさい。
  \\

  リスク分担型企業年金において、脱退一時金を繰り下げた場合の脱退一時金の額には、繰下
げを終了したときの調整率を適用する。
\end{truefalse}

\begin{answerbox}
  \textbf{正}
\end{answerbox}

\begin{shortanswer}
  確定給付企業年金（移換先確定給付企業年金）が、他の確定給付企業年金（移換元確定給付企
業年金）の加入者の資格を喪失した中途脱退者に係る脱退一時金相当額の移換を受けた場合で
あって、脱退一時金相当額の算定の基礎となった期間の一部を移換先確定給付企業年金の加入
者期間に算入する場合において、当該加入者期間に算入する期間の算定方法が満たすべき要件
として、確定給付企業年金法施行規則第８９条の４に定められている内容を３つ簡潔に入力し
なさい。（300 字以内）
\end{shortanswer}

\begin{answerbox}
  \begin{itemize}
    \item 移換先確定給付企業年金の規約に照らして当該移換された脱退一時金相当額の算定の基礎とな
る期間を算定すること。ただし、算定された期間が移換された脱退一時金相当額の算定の基礎
となった期間を超える場合にあっては、当該算定の基礎となった期間とすること。
    \item 脱退一時金相当額の算定の基礎となった期間を算入しないこととする場合にあっては、移換先
確定給付企業年金の加入者であった期間が１年未満である者に限り、その旨を規約で定めるこ
と。
    \item その他当該中途脱退者について不当に差別的なものでなく合理的な計算方法であると認められ
ること。
  \end{itemize}
\end{answerbox}

\newpage